\vspace{-1em}
\section{Referencias bibliográficas}

\begingroup
\renewcommand{\bibsection}{}
\begin{thebibliography}{99}
\bibitem[Fumeron et al.(2020)]{Fumeron2020} 
Fumeron, S., Berche, B., \& Moraes, F. (2020). Improving student understanding of electrodynamics: The case for differential forms. \textit{American Journal of Physics}, 88(12), 1083--1093. \url{https://doi.org/10.1119/10.0001754}

\bibitem[Griffiths(2017)]{Griffiths2017} 
Griffiths, D. J. (2017). \textit{Introduction to electrodynamics} (4.ª ed.). Cambridge University Press.

\bibitem[Hon \& Goldstein(2020)]{Hon2020} 
Hon, G., \& Goldstein, B. R. (2020). \textit{Reflections on the practice of physics: James Clerk Maxwell’s methodological odyssey in electromagnetism} (1.ª ed.). Routledge. \url{https://doi.org/10.4324/9780429351013}

\bibitem[Taflove \& Hagness(2005)]{Taflove2005} 
Taflove, A., \& Hagness, S. C. (2005). \textit{Computational electrodynamics: The finite-difference time-domain method} (3.ª ed.). Artech House.

\bibitem[Teixeira et al.(2023)]{Teixeira2023} 
Teixeira, F. L., Sarris, C., Zhang, Y., Na, D.-Y., Bérenger, J.-P., Su, Y., Okoniewski, M., Chew, W. C., Backman, V., \& Simpson, J. J. (2023). Finite-difference time-domain methods. \textit{Nature Reviews Methods Primers}, 3, Artículo 75. \url{https://doi.org/10.1038/s43586-023-00257-4}

\bibitem[Yee(1966)]{Yee1966} 
Yee, K. S. (1966). Numerical solution of initial boundary value problems involving Maxwell’s equations in isotropic media. \textit{IEEE Transactions on Antennas and Propagation}, 14(3), 302--307. \url{https://doi.org/10.1109/TAP.1966.1138693}

\bibitem[Balanis(2016)]{Balanis2016} 
Balanis, C. A. (2016). \textit{Antenna Theory: Analysis and Design} (4th ed.). Wiley.

\bibitem[International Telecommunication Union(2025)]{ITU2025} 
International Telecommunication Union, Radiocommunication Sector. (2025). \textit{Recommendation ITU-R P.2040-4: Effects of building materials and structures on radio-wave propagation in the range of 1 MHz to 450 GHz}. \url{https://www.itu.int/rec/R-REC-P.2040-4-202509-I/en}

\bibitem[Jackson(1998)]{Jackson1998} 
Jackson, J. D. (1998). \textit{Classical Electrodynamics} (3rd ed.). Wiley.

\bibitem[Mohammadkhani(2023)]{Mohammadkhani2023} 
Mohammadkhani, A. (2023). Axial ratio improvement of circularly polarized antenna using parasitic elements for 5G applications. \textit{International Journal of RF and Microwave Computer-Aided Engineering}, 2023, Artículo 6619073. \url{https://doi.org/10.1155/2023/6619073}

\bibitem[Philbin(2011)]{Philbin2011} 
Philbin, T. G. (2011). Electromagnetic energy–momentum in dispersive media. \textit{Physical Review A}, 83(1), Artículo 013823. \url{https://doi.org/10.1103/PhysRevA.83.013823}

\bibitem[Bhosale, Zhao \& Zhang(2024)]{Bhosale2024} 
Bhosale, A. A., Zhao, Y., \& Zhang, X. (2024). Electric field and SAR reduction in high-impedance RF arrays by using high permittivity materials for 7T MR imaging. \textit{PLOS ONE}, 19(7), Artículo e0305464. \url{https://doi.org/10.1371/journal.pone.0305464}

\bibitem[Diao, Liu \& Warnick(2023)]{Diao2023} 
Diao, J., Liu, L., \& Warnick, K. F. (2023). Receiving dipole antenna analysis using the Poynting streamline method: Providing new insights. \textit{IEEE Antennas and Propagation Magazine}. \url{https://doi.org/10.1109/MAP.2023.3260637}

\bibitem[MEEP Documentation(s.f.)a]{MeepBasics} 
MEEP Documentation. (s.f.). \textit{Python Tutorials/Basics}. Recuperado en 2026, de \url{https://meep.readthedocs.io/en/latest/Python_Tutorials/Basics/}

\bibitem[MEEP Documentation(s.f.)b]{MeepUI} 
MEEP Documentation. (s.f.). \textit{Python User Interface}. Recuperado en 2026, de \url{https://meep.readthedocs.io/en/latest/Python_User_Interface/}

\bibitem[Mörée \& Leijon(2022)]{Moree2022} 
Mörée, G., \& Leijon, M. (2022). Comparison of Poynting's vector and the power flow used in electrical engineering. \textit{AIP Advances}, 12(8), Artículo 085219. \url{https://doi.org/10.1063/5.0101339}

\bibitem[NanoComp/meep(2019)]{MeepIssue1276} 
NanoComp/meep. (2019). \textit{Absorbed power density tutorial} [Discusión técnica, Issue \#1276]. GitHub. \url{https://github.com/NanoComp/meep/issues/1276}

\bibitem[Oskooi et al.(2010)]{Oskooi2010} 
Oskooi, A. F., Roundy, D., Ibanescu, M., Bermel, P., Joannopoulos, J. D., \& Johnson, S. G. (2010). Meep: A flexible free-software package for electromagnetic simulations by the FDTD method. \textit{Computer Physics Communications}, 181(3), 687--702. \url{https://doi.org/10.1016/j.cpc.2009.11.008}

\bibitem[Poynting(1884)]{Poynting1884} 
Poynting, J. H. (1884). On the transfer of energy in the electromagnetic field. \textit{Philosophical Transactions of the Royal Society of London}, 175, 343--361. \url{https://doi.org/10.1098/rstl.1884.0016}

\bibitem[Roy, Niture \& Wu(2021)]{Roy2021} 
Roy, B., Niture, S., \& Wu, M. H. (2021). Biological effects of low power nonionizing radiation: A narrative review. \textit{Journal of Radiation Research and Imaging}, 1(1), 1--23.

capitulo 4
capitulo 5

\bibitem[Berenger(1994)]{Berenger1994} 
Berenger, J.-P. (1994). A perfectly matched layer for the absorption of electromagnetic waves. \textit{Journal of Computational Physics, 114}(2), 185--200. \url{https://doi.org/10.1006/jcph.1994.1159}

\bibitem[Kolar-Požun \& Kosec(2025)]{KolarPozun2025} 
Kolar-Po\v{z}un, A., \& Kosec, G. (2025). An RBF-based method for computational electromagnetics with reduced numerical dispersion. \textit{arXiv}. \url{https://arxiv.org/abs/2508.18205}

\bibitem[Roden \& Gedney(2000)]{Roden2000} 
Roden, J. A., \& Gedney, S. D. (2000). Convolutional PML (CPML): An efficient FDTD implementation of the CFS-PML for arbitrary media. \textit{Microwave and Optical Technology Letters, 27}(5), 334--338. \url{https://doi.org/10.1002/1098-2760(20001205)27:5<334::AID-MOP14>3.0.CO;2-A}

\bibitem[Schneider(2025)]{Schneider2025} 
Schneider, J. B. (2025). \textit{Understanding the Finite-Difference Time-Domain Method}. \url{https://eecs.wsu.edu/~schneidj/ufdtd/ufdtd.pdf}

\bibitem[MEEP Documentation(s.f.)c]{MeepCylindrical} 
MEEP Documentation. (s.f.). \textit{Cylindrical Coordinates}. Recuperado en 2026, de \url{https://meep.readthedocs.io/en/latest/Python_Tutorials/Cylindrical_Coordinates/}

\bibitem[MEEP Documentation(s.f.)d]{MeepSymmetry} 
MEEP Documentation. (s.f.). \textit{Exploiting Symmetry}. Recuperado en 2026, de \url{https://meep.readthedocs.io/en/latest/Exploiting_Symmetry/}

\bibitem[MEEP Documentation(s.f.)e]{MeepFAQ} 
MEEP Documentation. (s.f.). \textit{FAQ}. Recuperado en 2026, de \url{https://meep.readthedocs.io/en/latest/FAQ/}

\bibitem[MEEP Documentation(s.f.)f]{MeepMaterials} 
MEEP Documentation. (s.f.). \textit{Materials}. Recuperado en 2026, de \url{https://github.com/NanoComp/meep/blob/master/doc/docs/Materials.md}

\bibitem[MEEP Documentation(s.f.)g]{MeepParallel} 
MEEP Documentation. (s.f.). \textit{Parallel Meep}. Recuperado en 2026, de \url{https://meep.readthedocs.io/en/latest/Parallel_Meep/}

\bibitem[Oskooi et al.(2020)]{Oskooi2020Machine} 
Oskooi, A., Hogan, C., Hammond, A. M., Reid, M. T. H., \& Johnson, S. G. (2020). Factorized machine learning for performance modeling of massively parallel heterogeneous physical simulations. \textit{arXiv}. \url{https://arxiv.org/abs/2003.04287}

\bibitem[Fajardo et al.(2017)]{Fajardo2017} 
Fajardo, J. E., Vericat, F., Irastorza, G., Carlevaro, C. M., \& Irastorza, R. M. (2017). Sensitivity analysis on imaging the calcaneus using microwaves. \textit{arXiv}. \url{https://arxiv.org/abs/1709.04934}

\bibitem[Rashed et al.(2020)]{Rashed2020} 
Rashed, E. A., Diao, Y., \& Hirata, A. (2020). Learning-based estimation of dielectric properties and tissue density in head models for personalized radio-frequency dosimetry. \textit{Physics in Medicine and Biology, 65}(6), Artículo 065001. \url{https://doi.org/10.1088/1361-6560/ab6e53}

\bibitem[Hossain et al.(2026)]{hossain2026} Hossain, M. A., et al. (2026). A Simulation-Based Computational Study on the Dielectric Response of Human Hand Tissues to Radiofrequency Radiation from Mobile Devices. \textit{MDPI Sensors}, 26(1), 11.

\bibitem[Garcia et al.(2026)]{garcia2026} Garcia, L., et al. (2026). Impact of 5G mmWave radiation on human tissue using skin, cornea (eye) and enamel (tooth) as study candidates. \textit{ResearchGate}.

\bibitem[Wang et al.(2024)]{wang2024} Wang, H., et al. (2024). Permittivity Measurement in Multi-Phase Heterogeneous Concrete Using Evidential Regression Deep Network and High-Frequency Electromagnetic Waves. \textit{PMC}, 12387929.

\bibitem[Fichte et al.(2021)]{fichte2021} Fichte, C. O., et al. (2021). Dielectric material characterization of concrete in GHz range in dependence on pore volume and water content. \textit{Construction and Building Materials}, 313, 125442.

\bibitem[Smith et al.(2022)]{smith2022} Smith, A., et al. (2022). Design and analysis of wideband sub-6 GHz microstrip antenna. \textit{ResearchGate}.

\bibitem[Chen y Wang(2023)]{chen2023} Chen, Y., \& Wang, J. (2023). Poynting vector analysis of sub-6 GHz RF energy propagation in lossy biological tissues. \textit{IEEE Transactions on Antennas and Propagation}, 71(5), 4120-4131.

\bibitem[Jones et al.(2026)]{jones2026} Jones, R., et al. (2026). Sub-6-GHz 5G Large-Scale Path Loss Model for Shoemaker Rim F: Sensitivity to Transmitter Antenna Pattern. \textit{MDPI}, 7(1), 21.

\bibitem[Oskooi et al.(2023)]{oskooi2023} Oskooi, A., et al. (2023). Advances in the MEEP FDTD software package for complex electromagnetic simulations. \textit{Journal of Computational Physics}, 450, 110-125.

\end{thebibliography}
\endgroup

\noindent [2] S. T. Thornton y J. B. Marion, \textit{Classical dynamics of particles and systems}, 5ª ed. Belmont, CA, USA: Thomson Brooks/Cole, 2004.

\noindent [3] V. C. Barbosa y P. C. G. de Moraes, ``A Newtonian description of the motion of a spherical pendulum,'' \textit{Rev. Bras. Ens. Fís.}, vol. 32, no. 3, pp. 330--337, 2010. doi: \url{10.1590/S1806-11172010000300007}

\noindent [4] A. Hannachi, \textit{Patterns identification and data mining in weather and climate}. Cham, Suiza: Springer, 2021.

\noindent [5] V. A. Dobrushkin, \textit{MATHEMATICA tutorial for the second course in applied differential equations}, Center for Fluid Mechanics, Brown University. [En línea]. Disponible: \url{https://www.cfm.brown.edu/people/dobrush/am34/Mathematica/ch3/spherical.html}

\noindent [6] M. Meerstetter, \textit{A spherical pendulum: Modeling \& control}, Tesis de maestría, Lund Inst. of Technology, Lund Univ., Lund, Suecia, 2005. [En línea]. Disponible: \url{https://lup.lub.lu.se/luur/download?func=downloadFile&recordOId=8847877&fileOId=8859401}

\noindent [7] C. M. Bishop, \textit{Pattern recognition and machine learning}. New York, NY, USA: Springer, 2006.

\noindent [8] J. Brownlee, \textit{Deep learning for time series forecasting: Predict the future with MLPs, CNNs, and LSTMs in Python}. Vermont, Australia: Machine Learning Mastery, 2018.

\noindent [9] F. Aghware \textit{et al.}, ``Enhancing the Random Forest Model via Synthetic Minority Oversampling Technique for Credit-Card Fraud Detection,'' \textit{J. Comput. Theor. Appl.}, vol. 2, pp. 190--203, 2024. doi: \url{10.62411/jcta.10323}

\noindent [10] T. Chen y C. Guestrin, ``XGBoost: A scalable tree boosting system,'' en \textit{Proc. 22nd ACM SIGKDD Int. Conf. Knowl. Discov. Data Min.}, 2016, pp. 785--794. doi: \url{10.1145/2939672.2939785}

\noindent [11] J.-J. Liu y J.-C. Liu, ``Permeability Predictions for Tight Sandstone Reservoir Using Explainable Machine Learning and Particle Swarm Optimization,'' \textit{Geofluids}, vol. 2022, pp. 1--15, 2022. doi: \url{10.1155/2022/2263329}

\noindent [12] I. Goodfellow, Y. Bengio y A. Courville, \textit{Deep learning}. Cambridge, MA, USA: MIT Press, 2016. [En línea]. Disponible: \url{https://www.deeplearningbook.org}

\noindent [13] J. McGonagle, J. Alonso y S. Mollick, \textit{Feedforward Neural Networks | Brilliant Math \& Science Wiki}. [En línea]. Disponible: \url{https://brilliant.org/wiki/feedforward-neural-networks/}

\noindent [14] A. Habib y U. Yildirim, ``Developing a physics-informed and physics-penalized neural network model for preliminary design of multi-stage friction pendulum bearings,'' \textit{Eng. Appl. Artif. Intell.}, vol. 113, art. no. 104953, 2022. doi: \url{10.1016/j.engappai.2022.104953}

\noindent [15] N. Heimann, J. Petermann, D. Hartwig, R. Schnabel y L. Mathey, ``Predicting the motion of a high-Q pendulum subject to seismic perturbations using machine learning,'' \textit{Appl. Phys. Lett.}, vol. 122, no. 25, 2023. doi: \url{https://doi.org/10.48550/arXiv.2301.07694}

\noindent [16] G. M. Moatimid y T. S. Amer, ``Analytical approximate solutions of a magnetic spherical pendulum: Stability analysis,'' \textit{J. Vib. Eng. Technol.}, vol. 11, pp. 2155--2165, 2023. doi: \url{10.1007/s42417-022-00693-8}

\noindent [17] SciPy Community, \textit{scipy.integrate.odeint}, 2025. [En línea]. Disponible: \url{https://docs.scipy.org/doc/scipy/reference/generated/scipy.integrate.odeint.html}

\noindent [18] A. Hindmarsh, \textit{ODEPACK: Fortran ODE Solvers}, COMPUTING. [En línea]. Disponible: \url{https://computing.llnl.gov/projects/odepack}

\noindent [19] A. C. Hindmarsh, \textit{ODEPACK, a systematized collection of ODE solvers}, \textit{IMACS Trans. Sci. Comput.}, vol. 1, UCRL-88007, 1982. [En línea]. Disponible: \url{https://computing.llnl.gov/sites/default/files/ODEPACK_pub1_u88007.pdf}

\noindent [20] D. G. Zill y M. R. Cullen, \textit{Ecuaciones Diferenciales con problemas de valores en la frontera}, 9ª ed. Boston, MA, USA: Cengage Learning, 2014.
